\documentclass[11pt]{article}
\usepackage[T1]{fontenc}
\usepackage{lmodern}
\usepackage[margin=0.9in]{geometry}
\usepackage{amsmath,amsthm,amssymb}
\usepackage{microtype}
\usepackage{booktabs}
\usepackage{hyperref}
\hypersetup{colorlinks=true,linkcolor=black,citecolor=black,urlcolor=blue,
  pdftitle={Fibonacci layers generated by doubling and two reflections},
  pdfauthor={Mark Vladimir Prizmic}}
\setlength{\emergencystretch}{2em}
\newtheorem{theorem}{Theorem}
\newtheorem{lemma}[theorem]{Lemma}
\newtheorem{proposition}[theorem]{Proposition}
\newtheorem{corollary}[theorem]{Corollary}
\theoremstyle{remark}
\newtheorem{remark}[theorem]{Remark}
\newcommand{\Z}{\mathbb Z}
\newcommand{\ind}{\mathbf 1}

\title{Fibonacci layers generated by doubling\\and two reflections}
\author{Mark Vladimir Prizmic}
\date{November 27, 2026\quad\textnormal{\small Version 0.1}}

\begin{document}
\maketitle

\begin{abstract}
Starting at $1$, apply the maps $x\mapsto2x$, $x\mapsto1-x$, and
$x\mapsto2-x$, placing each integer in the first layer in which it is
reached. We give an explicit recursive formula for the shortest-path
distance from $1$ to any integer. The formula leads to the generating
function $1+2z/(1-z-z^2)$ for the layer sizes, proving that layer $n$ has
$2F_{n-1}$ elements for $n\ge2$. We also prove that exactly half of those
elements are positive and determine the even and odd counts, with the
initial exceptions stated explicitly. These conclusions establish the
layer-count and sign assertions recorded in OEIS A242364 and the parity
assertions from layer $4$ onward. The proof is self-contained and is
accompanied by reproducible finite verification code.
\end{abstract}

\section{The generation system}

For $x\in\Z$, put
\[
D(x)=2x,\qquad R(x)=1-x,\qquad S(x)=2-x.
\]
Define $G_1=\{1\}$ and, for $n\ge2$,
\begin{equation}\label{eq:layers}
G_n=\bigl(D(G_{n-1})\cup R(G_{n-1})\cup S(G_{n-1})\bigr)
       \setminus\bigcup_{i=1}^{n-1}G_i.
\end{equation}
Thus $G_n$ contains the integers first reached after $n-1$ operations.
The first layers are
\[
\begin{aligned}
G_1&=\{1\},&G_2&=\{0,2\},\\
G_3&=\{-1,4\},&G_4&=\{-3,-2,3,8\}.
\end{aligned}
\]
This is the generation system in Clark Kimberling's OEIS
A242364~\cite{oeis}. Equation~\eqref{eq:layers} makes the exclusion of
earlier layers explicit. That exclusion is implemented in the entry's
Mathematica program and agrees with its displayed rows; the prose
recurrence alone does not include the subtraction.

Let $F_0=0$, $F_1=1$, and $F_{m+1}=F_m+F_{m-1}$ for $m\ge1$.
Our main conclusion is the following.

\begin{theorem}\label{thm:main}
Every integer belongs to exactly one of the layers in
\eqref{eq:layers}. For every $n\ge2$,
\begin{equation}\label{eq:main-count}
|G_n|=2F_{n-1},\qquad
|G_n\cap\Z_{>0}|=F_{n-1}.
\end{equation}
For every $n\ge4$,
\begin{equation}\label{eq:main-parity}
|G_n\cap2\Z|=2F_{n-2},\qquad
|G_n\setminus2\Z|=2F_{n-3}.
\end{equation}
The even and odd counts in layers $1,2,3$ are respectively
$(0,1)$, $(2,0)$, and $(1,1)$.
\end{theorem}

In particular, exactly half the elements of each layer $n\ge2$ are
positive. In layer $2$, the other element is $0$; in every later layer,
the other half are negative. The restriction $n\ge4$ in
\eqref{eq:main-parity} is necessary.

\section{A recursive distance function}

For nonzero $x\in\Z$, define
\[
\varepsilon(x)=
\begin{cases}
1,&x=-2^t\text{ for some integer }t\ge0,\\
0,&\text{otherwise}.
\end{cases}
\]
Define $h:\Z\to\Z_{\ge0}$ recursively by
\begin{equation}\label{eq:h-base}
h(0)=1,\qquad h(1)=0,
\end{equation}
and, for $x\ne0$, by
\begin{equation}\label{eq:h-rec}
h(2x)=h(x)+1,\qquad
h(1-2x)=h(x)+2-\varepsilon(x).
\end{equation}
These equations determine $h$ at every integer. For a nonzero even
argument $y$, its recursive parent is $y/2$; for an odd argument $y\ne1$,
its parent is $(1-y)/2$. The parent's absolute value is smaller except
when $y=-1$, whose parent is the already defined $1$. Every increment
is positive. Thus the recursion is well-defined and nonnegative.

The exceptional odd arguments are precisely
\begin{equation}\label{eq:seeds}
\mathcal Q=\{2^u+1:u\ge1\}.
\end{equation}
The following values will be used repeatedly:
\begin{equation}\label{eq:explicit}
\begin{aligned}
h(2^t)&=t,& h(-2^t)&=t+2 &&(t\ge0),\\
h(2^u+1)&=u+2 &&&&(u\ge1).
\end{aligned}
\end{equation}
Indeed, $h(-1)=h(1)+2=2$, after which the first two formulas follow
by doubling. The last follows by taking the recursive parent
$-2^{u-1}$.

\begin{lemma}\label{lem:R}
For every $x\in\Z$,
\[
|h(1-x)-h(x)|\le1.
\]
\end{lemma}

\begin{proof}
If $x=2k\ne0$, equation~\eqref{eq:h-rec} gives
\[
h(1-2k)-h(2k)=1-\varepsilon(k)\in\{0,1\}.
\]
The pair $x=0,1$ is immediate from \eqref{eq:h-base}. Every other
case is the reverse of an even case, because $R$ is an involution.
\end{proof}

\begin{lemma}\label{lem:sign}
For every positive integer $k$,
\[
h(-k)-h(k)=2\quad\text{if }k\text{ is a power of }2.
\]
If $k$ is not a power of $2$, then
\[
|h(-k)-h(k)|\le1.
\]
\end{lemma}

\begin{proof}
The power-of-two assertion follows from \eqref{eq:explicit}.
Doubling preserves the difference $h(-k)-h(k)$, so in the other case
it suffices to consider an odd $k=2\ell+1\ge3$. Equation~\eqref{eq:h-rec}
gives
\begin{equation}\label{eq:sign-difference}
h(-k)-h(k)
=h(\ell+1)-h(-\ell)+\ind_{\{\ell\text{ is a power of }2\}}.
\end{equation}
If $\ell$ is not a power of $2$, Lemma~\ref{lem:R} applies to the
reflection pair $-\ell,\ell+1$ and proves the result.

If $\ell=1$, the right side of \eqref{eq:sign-difference} is
$h(2)-h(-1)+1=1-2+1=0$. Finally, if $\ell=2^t$ with $t\ge1$,
equation~\eqref{eq:explicit} gives
$h(\ell+1)=h(-\ell)=t+2$. The right side of
\eqref{eq:sign-difference} is then $1$.
\end{proof}

\begin{lemma}\label{lem:S}
For every $x\in\Z$,
\[
|h(2-x)-h(x)|\le1.
\]
\end{lemma}

\begin{proof}
First let $x=2k$ be even. If $k\notin\{0,1\}$, doubling reduces the
difference to
\[
h(2-2k)-h(2k)=h(1-k)-h(k),
\]
whose absolute value is at most $1$ by Lemma~\ref{lem:R}.
For $x=0$ or $2$, both heights are $1$.

The odd value $x=1$ is fixed by $S$. Every other odd reflection pair
has the form $2k+1,1-2k$ with $k\ge1$. For that pair,
\[
h(1-2k)-h(2k+1)
=h(k)-h(-k)+\ind_{\{k\text{ is a power of }2\}}.
\]
If $k$ is a power of $2$, the difference is $-1$ by
Lemma~\ref{lem:sign}. Otherwise its absolute value is at most $1$
by the same lemma. Reversing the pair covers negative odd $x$.
\end{proof}

\begin{proposition}\label{prop:distance}
Every integer $x$ can be reached from $1$ using $D,R,S$, and $h(x)$
is the minimum number of operations required.
\end{proposition}

\begin{proof}
On a $D$ edge, $h$ increases by $1$, except at the loop $0\mapsto0$.
On an $R$ or $S$ edge it increases by at most $1$, by
Lemmas~\ref{lem:R} and~\ref{lem:S}. Since $h(1)=0$, every path from
$1$ to $x$ therefore has length at least $h(x)$.

For the converse, we construct a path of length $h(x)$, using induction
on absolute value. The base paths are the empty path to $1$,
$1\xrightarrow{R}0$, and
$1\xrightarrow{D}2\xrightarrow{R}-1$.
For a nonzero even $x$ other than a base value, append $D$ to a path
to $x/2$. For an odd $x\ne1$ outside $\mathcal Q$, append $D$ and
then $R$ to a path to $(1-x)/2$. The recursion gives the required
lengths, and the parent has smaller absolute value.

For an exceptional $x=2^u+1\in\mathcal Q$, use $u$ doublings,
followed by $R$ and $S$:
\[
1\xrightarrow{D^u}2^u\xrightarrow{R}1-2^u
  \xrightarrow{S}1+2^u.
\]
Its length is $u+2=h(x)$ by \eqref{eq:explicit}. This proves
attainability and the claimed minimum.
\end{proof}

It follows that $G_n=\{x\in\Z:h(x)=n-1\}$. Every integer is reached,
and the earlier-layer exclusion makes its layer unique. Each level of
$h$ is finite: there are only finitely many words of bounded length in
the three maps. This also justifies the formal power series used below.

\section{The layer generating function}

Let
\[
C_d=|\{x\in\Z:h(x)=d\}|,\qquad
A(z)=\sum_{d\ge0}C_dz^d=\sum_{x\in\Z}z^{h(x)}.
\]
We decompose the integers into $1$, $0$, the nonzero even integers,
and the odd integers other than $1$.

The first two contributions are $1$ and $z$. Every nonzero even
integer is uniquely $2x$ with $x\ne0$, so its contribution is
$z(A(z)-z)$. Every odd integer other than $1$ is uniquely $1-2x$
with $x\ne0$. If every such child had height $h(x)+2$, this contribution
would be $z^2(A(z)-z)$.

There is exactly one exceptional parent $-2^t$ of height $t+2$ for
each $t\ge0$. Its odd child has height $t+3$, instead of the ordinary
height $t+4$. Hence the precise identity is
\begin{equation}\label{eq:gf-rec}
A(z)=1+z+z(A(z)-z)+z^2(A(z)-z)
            -\frac{z^4}{1-z}+\frac{z^3}{1-z}.
\end{equation}
All operations here are in the ring of formal power series; no analytic
convergence is required. The two correction terms have difference $z^3$,
and \eqref{eq:gf-rec} simplifies to
\[
(1-z-z^2)A(z)=1+z-z^2.
\]
Therefore
\begin{equation}\label{eq:gf}
\boxed{A(z)=1+\frac{2z}{1-z-z^2}.}
\end{equation}
Since $z/(1-z-z^2)=\sum_{d\ge1}F_dz^d$, we obtain
$C_0=1$ and $C_d=2F_d$ for $d\ge1$. In the original indexing,
$|G_n|=C_{n-1}=2F_{n-1}$ for $n\ge2$.

\section{Signs and parity}

Write
\[
P_d=|\{x>0:h(x)=d\}|,\qquad
M_d=|\{x<0:h(x)=d\}|.
\]
Every negative even integer is the double of a negative integer one
level earlier. Every negative odd integer is $1-2k$ for a unique
positive $k$, and has height $h(k)+2$, without an exceptional correction.
Consequently
\begin{equation}\label{eq:sign-count}
M_d=M_{d-1}+P_{d-2},
\end{equation}
where counts with negative subscripts are interpreted as zero.

The initial values are $P_0=P_1=1$ and $M_0=M_1=0$.
Thus $M_2=1$ and $M_3=2$. Equation~\eqref{eq:gf} and the fact that
$0$ occurs only at distance $1$ give
$P_2=M_2=F_2$ and $P_3=M_3=F_3$. Inductively, if $d\ge4$,
equation~\eqref{eq:sign-count} gives
\[
M_d=F_{d-1}+F_{d-2}=F_d,
\qquad P_d=C_d-M_d=F_d.
\]
The distance-$1$ layer $\{0,2\}$ also has exactly one positive element.
This proves the sign assertion in Theorem~\ref{thm:main}.

For $d\ge3$, the even integers of height $d$ are exactly the doubles
of all integers of height $d-1$. None of these parents is $0$, since
$h(0)=1$. Their number is therefore $C_{d-1}=2F_{d-1}$.
The number of odd integers of height $d$ is
\[
C_d-C_{d-1}=2F_d-2F_{d-1}=2F_{d-2}.
\]
Substituting $d=n-1$ proves \eqref{eq:main-parity} for $n\ge4$.
The first three layers displayed in Section~1 give the stated initial
exceptions. The proof of Theorem~\ref{thm:main} is complete.

\section{Verification and relation to earlier work}

\paragraph{Reproducible finite checks.}
The accompanying script \texttt{verify\_distance.py} uses only Python's
standard library. It computes the layers directly by breadth-first search,
using the three maps and a visited set. This generation step does not use
the proposed height function or Fibonacci counts. A separate recursive
evaluator of $h$ is then compared with every discovered distance.

The retained run covers distances $0$ through $25$, or layers $1$ through
$26$, and checks all $392\,835$ discovered integers. Every height,
layer count, positive count, and applicable parity count agrees. The
last three distance layers are:
\begin{center}
\begin{tabular}{rrrr}
\toprule
Distance $d$ & Total $C_d$ & Positive $P_d$ & Even \\
\midrule
23 & $57\,314$ & $28\,657$ & $35\,422$ \\
24 & $92\,736$ & $46\,368$ & $57\,314$ \\
25 & $150\,050$ & $75\,025$ & $92\,736$ \\
\bottomrule
\end{tabular}
\end{center}
The same run checks the three edge inequalities at every integer
$-100\,000\le x\le100\,000$. Its deterministic output is stored in
\texttt{verification.json}; it can be reproduced with
\begin{verbatim}
python verify_distance.py > verification.json
\end{verbatim}
These checks corroborate the implementation. The infinite statements
follow from the proofs above, not from a finite search.

\paragraph{Related work.}
Cox and McLellan~\cite{cox-mclellan} study the system generated by
$x\mapsto2x$ and $x\mapsto1-x$, without the second reflection
$x\mapsto2-x$. They give a recursive binary-tree description, prove
Fibonacci generation sizes for that two-map system, and determine its
generation index using binary Gray-code data. That published result
concerns a different system: the additional reflection changes the
distances considered here. The present proof is self-contained; its
exceptional odd children account for the added operation.

\paragraph{Source status and priority.}
The retrieved primary-source text of A242364 records the layer-count,
sign, and parity assertions as conjectures. The finite searches made
during preparation did not identify a proof of this exact three-map
statement. This is not an exhaustive literature review, and the
conjecture label alone does not establish novelty. This version makes
no claim of priority or of being the first proof.

\paragraph{AI assistance and review.}
The proof, manuscript, and verification code were developed with
substantial assistance from OpenAI's ChatGPT. Additional AI-assisted
reviews checked the distance recursion, boundary cases, and counting
argument. This version has not received external human peer review and
contains no proof-assistant certification. The full argument and code
are supplied for independent assessment.

\begin{thebibliography}{9}
\bibitem{oeis}
C.~Kimberling, \emph{A242364}, The On-Line Encyclopedia of Integer
Sequences, entry dated June 11, 2014.
\url{https://oeis.org/A242364}.
Primary indexed entry text consulted November 27, 2026; its program
and examples specify the removal of earlier generations.

\bibitem{cox-mclellan}
D.~Cox and K.~McLellan, \emph{A Problem on Generation Sets Containing
Fibonacci Numbers}, The Fibonacci Quarterly \textbf{55} (2017), no.~2,
105--113.
\url{https://www.fq.math.ca/Papers1/55-2/CoxMcLellan021717.pdf}.
\end{thebibliography}

\end{document}
